\begin{table}[htbp]
\centering
\caption{Phillips curve and Okun's law: stability across the six policy scenarios}
\label{tab:phillips_okun_scenarios}
\begin{threeparttable}
\begin{tabular}{llccc}
\toprule
Test & Scenario (rate / improvement policy) & $\hat{\beta}$ & Std.\ error & $p$-value \\
\midrule
Okun's law & High / Inactive & -0.182** & 0.081 & $0.0330$ \\
Okun's law & High / Active & -0.162** & 0.076 & $0.0440$ \\
Okun's law & Low / Inactive & -0.257*** & 0.083 & $0.0047$ \\
Okun's law & Low / Active & -0.209** & 0.078 & $0.0130$ \\
Okun's law & Medium / Inactive & -0.190*** & 0.067 & $0.0088$ \\
Okun's law & Medium / Active & -0.219** & 0.081 & $0.0123$ \\[4pt]
Phillips curve & High / Inactive & -0.00131*** & 0.00013 & $<0.001$ \\
Phillips curve & High / Active & -0.00118*** & 0.00015 & $<0.001$ \\
Phillips curve & Low / Inactive & -0.00117*** & 0.00018 & $<0.001$ \\
Phillips curve & Low / Active & -0.00128*** & 0.00017 & $<0.001$ \\
Phillips curve & Medium / Inactive & -0.00116*** & 0.00018 & $<0.001$ \\
Phillips curve & Medium / Active & -0.00126*** & 0.00015 & $<0.001$ \\
\bottomrule
\end{tabular}
\begin{tablenotes}
\footnotesize
\item Notes: same estimating equation as Table~\ref{tab:phillips_okun_main}, with the slope estimated separately within each of the six policy scenarios (regulated credit rate $\times$ improvement policy status) rather than pooled, as a stability check on the sign and significance of both stylised facts.
\item Every one of the twelve scenario-level estimates is negative and significant at $p<0.05$ (eight of the twelve at $p<0.01$), indicating the macroeconomic regularities targeted in calibration hold within each policy configuration individually, not only in the pooled panel.
\item *** $p<0.01$, ** $p<0.05$, * $p<0.10$.
\end{tablenotes}
\end{threeparttable}
\end{table}